
Predicting neural network accuracy directly from weights enables efficient model selection without running inference. This task requires handling permutation symmetry: hidden units can be reordered without changing network function. This work investigates whether permutation equivariance is merely helpful for data efficiency or categorically required for learning weight-to-accuracy mappings from raw weights.

Three methods are compared on a synthetic ResNet-20/CIFAR-10 model zoo of 6,000 models: a statistics baseline using per-layer aggregations with Ridge regression, an MLP baseline operating on flattened raw weights, and a permutation-equivariant Neural Functional Network (NFN) using DeepSets-style architecture. Experiments span training sizes N $\in$ {100, 250, 500, 1000, 2500, 5000} with 10 random seeds at the primary comparison point (N=500).

The results are categorical rather than quantitative. The MLP baseline achieves $R^2$ < 0 (worse than mean prediction) at all sample sizes including N=5000. The equivariant NFN achieves $R^2$ = 0.9985 at N=500 and maintains $R^2$ > 0.99 across all training sizes. The statistics baseline achieves $R^2$ = 0.9995 consistently. At N=500, the gap between NFN and MLP is $\Delta$ = 2.50 $R^2$ points (p = 4.58 $\times 10^{-6})$.

These findings reframe equivariance from a data efficiency technique to a fundamental requirement. Non-equivariant methods cannot learn weight-to-accuracy mappings from raw weights regardless of sample size within practical ranges.
